\chapter{Linear Structures and Abstract Data Types}

本章按数学四段体组织：对象--公理/不变式--定理--实例化。代数规约与物理实现分层，\texttt{Error}/\texttt{Exception}、数组下标、指针结构体均为实例化。

\section{栈 (Stack)}

\subsection{代数定义}
栈 $S$ 是定义在元素集 $E$ 上的代数结构，由以下算子与公理刻画：
\begin{itemize}
    \item \textbf{构造算子}：
    \begin{itemize}
        \item $\text{empty}: \to S$ （空栈）
        \item $\text{push}: S \times E \to S$ （压入）
    \end{itemize}
    \item \textbf{观测算子}：
    \begin{itemize}
        \item $\text{pop}: S \to S \cup \{\bot\}$ （弹出，偏函数）
        \item $\text{top}: S \to E \cup \{\bot\}$ （栈顶，偏函数）
        \item $\text{isEmpty}: S \to \mathbb{B}$ （判空）
    \end{itemize}
\end{itemize}

\textbf{公理系统}（$\forall s \in S, e \in E$）：
\begin{enumerate}
    \item $\text{isEmpty}(\text{empty}) = \text{true}$
    \item $\text{isEmpty}(\text{push}(s, e)) = \text{false}$
    \item $\text{pop}(\text{push}(s, e)) = s$
    \item $\text{top}(\text{push}(s, e)) = e$
    \item $\text{pop}(\text{empty}) = \bot, \quad \text{top}(\text{empty}) = \bot$
\end{enumerate}

\begin{itemize}
    \item \textbf{实例化}：网页 \texttt{stack} 演示 \texttt{push}/\texttt{pop}/\texttt{top} 的帧序列；语言实例中 $\bot$ 实现为异常或错误码，数学层仅以偏函数表达。
\end{itemize}

\subsection{物理实现方案}
\begin{enumerate}
    \item \textbf{顺序栈（有界实例）}：$S_m$ 映射到 $A[0..m-1]$，$t\in\{-1,\dots,m-1\}$，不变式 $t<m-1$ 为 \texttt{push} 的定义域，$t\ge0$ 为 \texttt{pop} 的定义域。
    \item \textbf{链栈（无界实例）}：映射到单向链表，表头即栈顶，无容量上限，为 $S$ 的无界实例。
\end{enumerate}

\begin{algorithm}[H]
\caption{栈操作规约（有界顺序实例，数学层为偏函数）}
\begin{algorithmic}[1]
\State \textbf{状态}: $A[0..m-1]$, $t\gets-1$
\Procedure{Push}{$e$} \Comment{requires $t<m-1$}
    \State $t\gets t+1$; $A[t]\gets e$
\EndProcedure
\Procedure{Pop}{} \Comment{requires $t\ge0$}
    \State $e\gets A[t]$; $t\gets t-1$; \Return $e$
\EndProcedure
\end{algorithmic}
\end{algorithm}

\begin{itemize}
    \item \textbf{实例化注}：溢出/下溢 $t=m-1$/$t=-1$ 时数学上无定义（$\bot$）；语言实例以 \texttt{Error}、异常或返回值标记，网页 \texttt{stack} 即按此演示。
\end{itemize}

\begin{figure}[H]
\centering
\begin{tikzpicture}[
    cell/.style={draw, minimum width=1.4cm, minimum height=0.75cm, font=\ttfamily},
    arrow/.style={->, >=latex, thick},
    lbl/.style={font=\small}
]
    \node[cell] (a) at (0,0) {$d_1$};
    \node[cell] (b) at (0,0.75) {$d_2$};
    \node[cell] (c) at (0,1.5) {$d_3$};
    \node[cell, draw=none, minimum height=0.4cm] (d) at (0,2.25) {};
    \draw (d.south west) -- (d.north west);
    \draw (d.south east) -- (d.north east);
    \draw[arrow] (-2.2, 2.25) -- node[lbl, below] {$\text{push}(e)$} (d.west);
    \draw[arrow] (d.east) -- node[lbl, below] {$\text{pop}()$} (2.2, 2.25);
    \node[lbl, anchor=west] at ($(c.east)+(0.2,0)$) {$\leftarrow$ 栈顶 ($t$)};
    \node[lbl, anchor=west] at ($(a.east)+(0.2,0)$) {$\leftarrow$ 栈底};
\end{tikzpicture}
\caption{栈 (Stack) 逻辑模型}
\label{fig:stack_model}
\end{figure}

\section{队列 (Queue)}

\subsection{代数定义}
队列 $Q$ 是定义在 $E$ 上的代数结构：
\begin{itemize}
    \item \textbf{构造算子}：$\text{emptyQ}: \to Q$，$\text{enqueue}: Q \times E \to Q$
    \item \textbf{观测算子}：$\text{dequeue}: Q \to Q \cup \{\bot\}$，$\text{front}: Q \to E \cup \{\bot\}$，$\text{isEmptyQ}: Q \to \mathbb{B}$（均为偏函数，空队列上为 $\bot$）
\end{itemize}

\textbf{公理系统}（$\forall q \in Q, e \in E$）：
\begin{enumerate}
    \item $\text{isEmptyQ}(\text{emptyQ}) = \text{true}$
    \item $\text{dequeue}(\text{enqueue}(\text{emptyQ}, e)) = \text{emptyQ}$
    \item $\text{dequeue}(\text{enqueue}(\text{enqueue}(q, e_1), e_2)) = \text{enqueue}(\text{dequeue}(\text{enqueue}(q, e_1)), e_2)$
    \item $\text{front}(\text{enqueue}(\text{emptyQ}, e)) = e$
    \item 若 $\text{isEmptyQ}(q) = \text{false}$，则 $\text{front}(\text{enqueue}(q, e)) = \text{front}(q)$
\end{enumerate}

\subsection{循环队列 (Circular Queue)}
为消除假溢出，将数组索引映射到剩余类环 $\mathbb{Z}_m$。
\begin{itemize}
    \item \textbf{对象}：$Q[0..m-1]$, $f,r\in\mathbb{Z}_m$。
    \item \textbf{不变式}：$f\leftarrow(f+1)\bmod m$, $r\leftarrow(r+1)\bmod m$；队空 $r=f$，队满 $(r+1)\bmod m=f$（牺牲一槽位的不变式实例，其他实例以计数器区分）。
    \item \textbf{定理}：环上 \texttt{enqueue}/\texttt{dequeue} 均为 $O(1)$。
    \item \textbf{实例化}：网页 \texttt{queue} 演示普通与循环队列的指针帧；计数器方案为另一实例。
\end{itemize}

\begin{algorithm}[H]
\caption{循环队列规约（有界实例，数学层为偏函数）}
\begin{algorithmic}[1]
\State \textbf{状态}: $Q[0..m-1]$, $f\gets0, r\gets0$
\Procedure{Enqueue}{$e$} \Comment{requires $(r+1)\bmod m\neq f$}
    \State $Q[r]\gets e$; $r\gets(r+1)\bmod m$
\EndProcedure
\Procedure{Dequeue}{} \Comment{requires $r\neq f$}
    \State $e\gets Q[f]$; $f\gets(f+1)\bmod m$; \Return $e$
\EndProcedure
\end{algorithmic}
\end{algorithm}

\begin{itemize}
    \item \textbf{实例化注}：满/空判据的 \texttt{Error} 为语言实例的 $\bot$ 实现；数学层以定义域表达。
\end{itemize}

\begin{figure}[H]
\centering
\begin{tikzpicture}[
    cell/.style={draw, minimum width=1.1cm, minimum height=0.75cm, font=\ttfamily},
    arrow/.style={->, >=latex, thick},
    lbl/.style={font=\small}
]
    \matrix (q) [matrix of nodes, nodes={cell}, column sep=-\pgflinewidth] {
        $d_1$ & $d_2$ & $d_3$ & $d_4$ & \phantom{A} & \phantom{A} \\
    };
    \draw[arrow] (q-1-1.west) -- ++(-4,0) node[lbl, midway, above] {$\text{dequeue}()$};
    \draw[arrow] ($(q-1-6.east)+(4,0)$) -- ++(-4,0) node[lbl, midway, above] {$\text{enqueue}(e)$};
    \node[lbl] (front) at ($(q-1-1.south)+(0,-0.9)$) {$f$};
    \draw[arrow] (front.north) -- (q-1-1.south);
    \node[lbl] (rear) at ($(q-1-4.south)+(0,-0.9)$) {$r$};
    \draw[arrow] (rear.north) -- (q-1-4.south);
\end{tikzpicture}
\caption{普通队列 (Queue)}
\label{fig:queue_model}
\end{figure}

\begin{figure}[H]
\centering
\begin{tikzpicture}[
    cell/.style={draw, minimum size=0.8cm, font=\ttfamily},
    arrow/.style={->, >=latex},
    lbl/.style={font=\small}
]
    \node[cell] (s5) at (90:1.8) {5};
    \node[cell] (s0) at (30:1.8) {0};
    \node[cell] (s1) at (-30:1.8) {1};
    \node[cell] (s2) at (-90:1.8) {2};
    \node[cell] (s3) at (-150:1.8) {3};
    \node[cell] (s4) at (150:1.8) {4};
    \draw[arrow] (s5) to[bend left=15] (s0);
    \draw[arrow] (s0) to[bend left=15] (s1);
    \draw[arrow] (s1) to[bend left=15] (s2);
    \draw[arrow] (s2) to[bend left=15] (s3);
    \draw[arrow] (s3) to[bend left=15] (s4);
    \draw[arrow] (s4) to[bend left=15] (s5);
    \node[lbl, anchor=west] at ($(s0.east)+(0.3,0)$) {$\leftarrow f$};
    \node[lbl, anchor=east] at ($(s3.west)+(-0.3,0)$) {$r \rightarrow$};
    \node[lbl, align=left, anchor=north west] at (3.2, -0.8) {
        队满判据 ($m=6$):\\
        $(r + 1) \bmod 6 = f$\\
        $(4 + 1) \bmod 6 = 5$
    };
\end{tikzpicture}
\caption{循环队列 (Circular Queue)}
\label{fig:queue_circular}
\end{figure}

\section{字符串 (String)}

\subsection{形式化定义}
\begin{itemize}
    \item \textbf{对象}：字母表 $\Sigma$ 有限非空。
    \item \textbf{字符串幺半群}：$\Sigma^* = \bigcup_{k=0}^{\infty} \Sigma^k$，空串 $\varepsilon\in\Sigma^0$，连接 $\circ: \Sigma^*\times\Sigma^*\to\Sigma^*$ 结合，$\varepsilon$ 为单位元。
    \item \textbf{不变式}：$|s\circ t| = |s|+|t|$，$|\varepsilon|=0$。
\end{itemize}

\begin{itemize}
    \item \textbf{物理实例}：底层以字符数组存储，边界以哨兵 $\#\notin\Sigma$ 界定。语言实例取 $\#=\backslash0$（C 字符串），数学层仅以 $\#\notin\Sigma$ 表达；网页 \texttt{string-ops} 演示连接、子串等幺半群操作。
\end{itemize}

\begin{figure}[H]
\centering
\begin{tikzpicture}[
    cell/.style={draw, minimum width=1.5cm, minimum height=0.75cm, font=\ttfamily},
    lbl/.style={font=\small}
]
    \matrix (str) [matrix of nodes, nodes={cell}, column sep=-\pgflinewidth] {
        $c_1$ & $c_2$ & $c_3$ & $\#$ \\
    };
    \node[lbl] at ($(str-1-1.south)+(0,-0.35)$) {0};
    \node[lbl] at ($(str-1-2.south)+(0,-0.35)$) {1};
    \node[lbl] at ($(str-1-3.south)+(0,-0.35)$) {2};
    \node[lbl] at ($(str-1-4.south)+(0,-0.35)$) {3};
    \node[lbl, anchor=west] at ($(str.east)+(0.3,0)$) {$\leftarrow s \in \Sigma^*$ 的物理映射（$\#\notin\Sigma$）};
\end{tikzpicture}
\caption{字符串物理存储布局}
\label{fig:string_layout}
\end{figure}

\section{高级映射与哈希表体系}

\subsection{字典 (Dictionary) / 映射 (Map)}
\begin{itemize}
    \item \textbf{对象}：偏函数 $\mathcal{M}: K \rightharpoonup V$，$\text{dom}(\mathcal{M})\subseteq K$。
    \item \textbf{不变式}：单射于键（$\mathcal{M}(k_1)=\mathcal{M}(k_2)\implies k_1=k_2$ 在值可逆时，或键唯一性由表保证）。
\end{itemize}

\subsection{哈希表 (Hash Table)}
通过 $h: K \to \mathbb{Z}_m$ 将键空间压缩至槽位 $\{0,\dots,m-1\}$。

\subsubsection{(1) 哈希函数}
\begin{itemize}
    \item \textbf{除留余数法}：$h(k)=k\bmod m$，$m$ 取质数以最小化冲突为实例化启发式。
    \item \textbf{理想哈希}：$\forall k_i\neq k_j\implies h(k_i)\neq h(k_j)$（鸽巢原理下仅当 $|K|\le m$ 可实现）。
\end{itemize}

\subsection{哈希碰撞 (Hash Collision)}
\subsubsection{(1) 碰撞必然性}
由鸽巢原理，若 $|K|>m$，则 $\exists k_i\neq k_j$ 使 $h(k_i)=h(k_j)$。

\begin{figure}[H]
\centering
\begin{tikzpicture}[
    keybox/.style={draw, minimum width=2.2cm, minimum height=0.7cm, font=\ttfamily},
    cell/.style={draw, minimum width=2cm, minimum height=0.7cm, font=\small},
    arrow/.style={->, >=latex, thick},
    lbl/.style={font=\small}
]
    \node[keybox] (k1) at (0, 0.55) {$k_1$};
    \node[keybox] (k2) at (0, -0.55) {$k_2$};
    \node[lbl, anchor=south] at ($(k1.north)+(0,0.15)$) {【输入键 $K$】};
    \node[cell] (s0) at (6.5, 0.7) {0: $\emptyset$};
    \node[cell, fill=yellow!25] (s1) at (6.5, 0) {1: $v$};
    \node[cell] (s2) at (6.5, -0.7) {2: $\emptyset$};
    \node[lbl, anchor=south] at ($(s0.north)+(0,0.3)$) {【槽位 $\mathbb{Z}_m$】};
    \draw[arrow] (k1.east) -- node[lbl, above, midway] {$h(k) \bmod m$} (s1.west);
    \draw[arrow] (k2.east) -- (s1.west);
    \node[lbl, red, anchor=west] at ($(s1.east)+(0.2,0)$) {$\leftarrow$ 碰撞：$h(k_1) = h(k_2)$};
\end{tikzpicture}
\caption{哈希碰撞三大核心要素模型}
\label{fig:hash_collision}
\end{figure}

\subsubsection{(2) 冲突解决方案}
\begin{enumerate}
    \item \textbf{开放寻址法 (Open Addressing)}
    \begin{itemize}
        \item 探测序列：$h_i(k) = (h(k) + \delta(i)) \bmod m$，所有槽位在主表中，冲突时按序列探测。
        \item \textbf{线性探测} $\delta(i)=i$：$h(k),h(k)+1,\dots$ 期望 $O(1+\alpha)$，但易主聚集。
        \item \textbf{二次探测} $\delta(i)=i^2$：$h(k),h(k)+1,h(k)+4,\dots$ 缓解主聚集；$m$ 为质数且 $\alpha<0.5$ 时可遍历。
        \item \textbf{双重散列} $h_i(k)=(h_1(k)+i\cdot h_2(k))\bmod m$，$h_2(k)\not\equiv0$，各键独立序列，最接近随机。
        \item \textbf{公共缺陷}：删除需墓碑标记，再散列为扩容实例。
        \item \textbf{实例化}：网页 \texttt{hash-table} 演示开放寻址的探测链；三者均为 $h_i$ 的实例化。
    \end{itemize}

    \item \textbf{链地址法 (Chaining)}
    \begin{itemize}
        \item 槽位 $i\in\mathbb{Z}_m$ 维护 $L_i=\{(k,v)\mid h(k)=i\}$。
        \item 期望 $O(1+\alpha)$，$\alpha=n/m$。
        \item \textbf{动态重构}：$|L_i|>\tau$ 时重构为平衡树（实例取红黑树），最坏 $O(\log n)$。
    \end{itemize}
\end{enumerate}

\begin{algorithm}[H]
\caption{哈希表查找与插入规约（以 $h_i$ 泛化，线性为实例）}
\begin{algorithmic}[1]
\State \textbf{状态}: $T[0..m-1]$, $h_i(k)$
\Procedure{Search}{$k$}
    \State $i\gets0$
    \Repeat
        \State $j\gets h_i(k)$ \Comment{实例：$(h(k)+i)\bmod m$ 为线性}
        \If{$T[j]=\emptyset$ \textbf{or} $T[j].key=k$} \Return $T[j]$ \EndIf
        \State $i\gets i+1$
    \Until{$i=m$}
    \State \Return $\bot$
\EndProcedure
\Procedure{Insert}{$k,v$}
    \State $item\gets\text{Search}(k)$
    \If{$item\neq\bot$} $item.value\gets v$
    \Else
        \State $i\gets0$; \While{True} $j\gets h_i(k)$; \If{$T[j]=\emptyset$} $T[j]\gets\langle k,v\rangle$; \Return \EndIf; $i\gets i+1$ \EndWhile
    \EndIf
\EndProcedure
\end{algorithmic}
\end{algorithm}

\begin{figure}[H]
\centering
\begin{tikzpicture}[
    cell/.style={draw, minimum width=1.1cm, minimum height=0.7cm, font=\ttfamily},
    kv/.style={draw, minimum width=1.9cm, minimum height=0.6cm, font=\small},
    arrow/.style={->, >=latex, thick},
    lbl/.style={font=\small}
]
    \matrix (bkt) [matrix of nodes, nodes={cell}, row sep=-\pgflinewidth] {
        0 \\
        1 \\
        2 \\
    };
    \node[lbl, anchor=south] at ($(bkt-1-1.north)+(0,0.15)$) {槽位 $i$};
    \node[kv, right=1.2cm of bkt-1-1] (ka) {$(k_1, v_1)$};
    \node[kv, right=0.5cm of ka] (kb) {$(k_2, v_2)$};
    \draw[arrow] (bkt-1-1.east) -- (ka.west);
    \draw[arrow] (ka.east) -- (kb.west);
    \node[lbl, anchor=west] at ($(kb.east)+(0.3,0)$) {$\to \emptyset$};
    \draw[arrow] (bkt-2-1.east) -- ++(1.2,0) node[lbl, anchor=west] {$\emptyset$};
    \node[draw, minimum width=2.2cm, minimum height=0.65cm, font=\small, right=1.2cm of bkt-3-1] (rbt) {红黑树};
    \draw[arrow] (bkt-3-1.east) -- (rbt.west);
    \node[lbl, anchor=north, align=center] at ($(rbt.south)+(0,-0.2)$) {
        （冲突数 $> \tau$ 时重构，\\
        检索复杂度 $O(\log n)$）
    };
\end{tikzpicture}
\caption{链地址法与动态重构模型}
\label{fig:hash_chaining}
\end{figure}

\begin{itemize}
    \item \textbf{实例化}：网页 \texttt{hash-table} 演示链地址与重构；三探测、墓碑、再散列均为本规约的实例化。
\end{itemize}
